\documentclass[12pt]{article}
\usepackage[utf8]{inputenc}
\usepackage[T1]{fontenc}
\usepackage[margin=1in]{geometry}
\usepackage{setspace}
\usepackage{xcolor}
\usepackage[colorlinks=true, urlcolor=blue!60!black, linkcolor=blue!60!black]{hyperref}
\setlength{\parindent}{2em}
\setlength{\parskip}{1.25em}
\renewcommand{\baselinestretch}{1.0}
\begin{document}
\begin{center}
{\LARGE \textit{NAACP v. Alabama ex rel. Patterson}, 357 U.S. 449 (1958)\par}
\vspace{0.8em}
{\large Case Brief --- Four-Page Edition\par}
\end{center}
\vspace{1em}

\section{Memory Jogger}

Alabama tried to oust the NAACP and held it in contempt for refusing to hand over its rank-and-file membership lists; a unanimous Court held that compelled disclosure of members' identities, exposing them to reprisal without a sufficient justification supported by a substantial relation between disclosure and an important state interest, violates the freedom of association protected by the Fourteenth Amendment (\href{https://supreme.justia.com/cases/federal/us/357/449/}{Justia}).

\section{Detailed Case Facts}

Alabama required out-of-state corporations to file their charters and designate an agent before doing intrastate business. The NAACP, a New York nonprofit membership corporation, had operated in Alabama through local affiliates since 1918 and opened a regional office in 1951 without qualifying, believing itself exempt.

\section{Procedural History}

In 1956 Attorney General John Patterson sued in Montgomery County Circuit Court to oust the NAACP, citing its support for Black students seeking admission to the University of Alabama and for the Montgomery bus boycott. The court issued an \textit{ex parte} order barring the NAACP from operating and from qualifying. The State then sought records including the names and addresses of all Alabama ``members'' and ``agents.'' The NAACP offered to qualify and produced its charter, business records, officers' names, total membership, and dues, but refused to disclose its rank-and-file members. Held in civil contempt, it was fined \$10,000, rising to \$100,000 after five days, and barred from a merits hearing until it purged the contempt. The Alabama Supreme Court twice refused review, saying mandamus was the proper remedy. The U.S. Supreme Court reversed on June 30, 1958.

\section{Judicial Votes}

Reversed and remanded, \textbf{9--0}.
\begin{itemize}
\item \textbf{Majority (9):} Justice Harlan, joined by Chief Justice Warren and Justices Black, Frankfurter, Douglas, Burton, Clark, Brennan, and Whittaker.
\item \textbf{Concurrences and dissents:} None.
\end{itemize}

\section{Holding}

The Alabama Supreme Court's novel procedural bar was not an adequate state ground, and the NAACP could assert its members' rights. Freedom to associate for the advancement of beliefs and ideas is part of the ``liberty'' protected by the Fourteenth Amendment's Due Process Clause. Compelled disclosure of rank-and-file lists, which exposed members to serious private reprisal, was unconstitutional because Alabama showed no compelling interest substantially related to the names. The contempt judgment and \$100,000 fine fell.

\section{Analysis of Opinions}

\textbf{Majority (Harlan).} \emph{State ground.} Alabama had long reviewed contempt orders by certiorari and had even told the NAACP that certiorari was the ``established'' method; ``[n]ovelty in procedural requirements cannot be permitted to thwart review in this Court.'' \emph{Standing.} The NAACP could assert members' anonymity because requiring them to sue in their own names would destroy the very right claimed, and the Association and its members were ``in every practical sense, identical.'' \emph{Association.} Given the close nexus between speech and assembly, state action curtailing association faces ``the closest scrutiny,'' and even indirect deterrence counts. Uncontroverted evidence showed that past disclosure exposed members to ``economic reprisal, loss of employment, threat of physical coercion, and other manifestations of public hostility.'' That reprisals came from private citizens did not matter, because private action takes hold only after state compulsion. \emph{Interest and fit.} Any justifying interest must be ``compelling.'' Alabama's purpose---deciding whether the NAACP did intrastate business---had no ``substantial bearing'' on members' names, since the NAACP had admitted its presence, offered to qualify, and produced everything else. \textit{Bryant v. Zimmerman}, upholding disclosure of Ku Klux Klan rosters, was distinguished by the Klan's violence and refusal to provide any information.

\textbf{Separate opinions.} None.

\textbf{Later developments.} \textit{Bates v. City of Little Rock} and \textit{Shelton v. Tucker} (1960) struck similar disclosure laws, while \textit{Buckley v. Valeo} and \textit{Doe v. Reed} upheld general disclosure under ``exacting scrutiny,'' with as-applied exemptions. \textit{Americans for Prosperity Foundation v. Bonta} (2021) required narrow tailoring even for confidential donor filings, and \textit{First Choice Women's Resource Centers v. Davenport} (2026) held that a donor subpoena itself inflicts a present associational injury.

\section{Examples: Future Applications}

\begin{enumerate}
\item \textbf{Same-Side: Legislative Subpoena of Members.} A committee subpoenas a rural LGBTQ advocacy group's full membership list after it has registered and produced its officers and budget, and members have faced job loss and threats. As in \textit{NAACP v. Alabama}, the names bear no substantial relation to the purpose.
\item \textbf{Same-Side: Teachers Listing Every Association.} A State requires teachers to list every organization joined in five years. \textit{Shelton v. Tucker} struck such a law as far broader than the interest in fitness.
\item \textbf{Opposite-Side: Federal Candidate Contributions.} A statute requires public reporting of contributors over \$200. Under \textit{Buckley}, disclosure survives exacting scrutiny, though a group showing likely reprisals may seek an exemption.
\item \textbf{Opposite-Side: Referendum Petition Signers.} State law makes signed petitions public. Under \textit{Doe v. Reed}, a facial challenge fails because signing is a public lawmaking act.
\item \textbf{Fence-Sitter: Civic Transparency Act.} A State requires nonprofits spending over \$5,000 on issue ads to file confidential lists of donors over \$250 to combat foreign influence, and subpoenas an environmental group. \textit{Davenport} supports standing and \textit{Bonta} condemns a dragnet; the State invokes \textit{Buckley}, \textit{Reed}, and confidentiality.
\end{enumerate}

\section{Critique}

\textbf{An unenumerated right.} ``Association'' does not appear in the First Amendment. Harlan derived it functionally, as an amplifier of speech and assembly anchored in Fourteenth Amendment ``liberty,'' laying groundwork later invoked in privacy cases such as \textit{Griswold v. Connecticut}.

\textbf{Shield or sword?} Critics, including Justice Sotomayor dissenting in \textit{Bonta}, argue that \textit{NAACP v. Alabama} rested on an extraordinary record of Jim Crow violence and pretextual targeting, and that applying facial exacting scrutiny to routine donor disclosure risks insulating ``dark money'' from legitimate transparency and anti-fraud oversight.

\section{Key Quotations}

\begin{quote}
``It is beyond debate that freedom to engage in association for the advancement of beliefs and ideas is an inseparable aspect of the `liberty' assured by the Due Process Clause of the Fourteenth Amendment, which embraces freedom of speech. \dots\ [S]tate action which may have the effect of curtailing the freedom to associate is subject to the closest scrutiny.''
\end{quote}
\noindent --- Justice Harlan, 357 U.S. at 460--461.

\begin{quote}
``Inviolability of privacy in group association may in many circumstances be indispensable to preservation of freedom of association, particularly where a group espouses dissident beliefs.''
\end{quote}
\noindent --- Justice Harlan, 357 U.S. at 462.

\begin{quote}
``To require that it be claimed by the members themselves would result in nullification of the right at the very moment of its assertion.''
\end{quote}
\noindent --- Justice Harlan, 357 U.S. at 459.

\end{document}
